本笔记是南京理工大学《常微分方程》课程的课堂笔记，由课堂讲义 PDF 整理而成。笔记以王高雄等《常微分方程》（第四版）为主线，逐章整理了常微分方程的基本理论与求解方法：第一章绪论介绍微分方程的背景与基本概念；第二章系统讲解一阶微分方程的初等积分法（变量分离、线性方程、恰当方程、隐方程等）；第三章讨论一阶微分方程解的存在唯一性、解的延拓与奇解；第四章研究高阶线性微分方程的一般理论与常系数方程的解法；第五章将线性理论推广到微分方程组，并介绍矩阵指数法与拉普拉斯变换；第六章概述非线性微分方程的稳定性与定性理论。笔记中的公式按数学含义重建，模糊处标有 $[?]$ 并附校对说明，供后续对照原讲义核实。

\section*{课程信息}

\begin{itemize}
    \item \textbf{教材}：王高雄等《常微分方程》（第四版），高等教育出版社
    \item \textbf{参考书}：
    \begin{enumerate}
        \item 朱思铭《常微分方程学习辅导与习题解答》
        \item 丁同仁《常微分方程教程》
        \item 王柔怀、伍卓群《常微分方程讲义》
        \item 庞特里亚金《常微分方程》（倪明康译）
        \item 张芷芬等《微分方程定性理论》
    \end{enumerate}
    \item \textbf{学时}：第 5--14 周，共 48 学时
    \item \textbf{考核}：总成绩 $=$ 平时 $20\%$ + 期中 $10\%$ + 期末 $70\%$；每周二交作业
    \item \textbf{课程内容}：第一章 绪论 $\rightarrow$ 第二章 一阶微分方程的初等积分法 $\rightarrow$ 第三章 解的存在定理 $\rightarrow$ 第四章 高阶微分方程 $\rightarrow$ 第五章 线性微分方程组 $\rightarrow$ 第六章 非线性微分方程 $\rightarrow$ 第七章 一阶线性偏微分方程$^{*}$ $\rightarrow$ 第八章 边值问题
    \item \textbf{课程目的}：掌握各类可求解方程（组）的类型与解法；熟悉解的基本性质（存在性、唯一性）；了解稳定性分析、定性分析等基本方法
\end{itemize}

\noindent 笔记由课堂 PDF 整理，公式按数学含义重建；模糊处标 $[?]$ 待校对。由于本课程主要研究常微分方程，除特别说明外，笔记中"微分方程"均指常微分方程。

待续事项：
常见模型的绘图。
